\documentclass{article}
\usepackage[margin=2cm]{geometry}
\usepackage{amsmath}
\usepackage{graphicx}
\usepackage{booktabs}
\usepackage[utf8]{inputenc}
\begin{document}
\(e_{i} \pm e_{j}\) for all \(i \neq j\). The free energy in the specific case of a vector multiplet in the adjoint, a single antisymmetric hypermultiplet, and \(N_{f}\) fundamental hypermultiplets then is
\[
\begin{aligned}
F\left(\lambda_{i}\right)= & \sum_{i \neq j}\left[F_{V}\left(\lambda_{i}-\lambda_{j}\right)+F_{V}\left(\lambda_{i}+\lambda_{j}\right)+F_{H}\left(\lambda_{i}-\lambda_{j}\right)+F_{H}\left(\lambda_{i}+\lambda_{j}\right)\right] \\
& +\sum_{i}\left[F_{V}\left(2 \lambda_{i}\right)+F_{V}\left(-2 \lambda_{i}\right)+N_{f} F_{H}\left(\lambda_{i}\right)+N_{f} F_{H}\left(-\lambda_{i}\right)\right]
\end{aligned}
\]
The next step is to look for extrema of this function in the specific case of \(\lambda_{i} \geq 0\) for all \(i\). Extrema in the case of non-positive \(\lambda_{i}\) can be obtained from these through action of the Weyl group. To calculate the extrema, one first assumes that as \(N \rightarrow \infty\), the vess scale as \(\lambda_{i}=N^{\alpha} x_{i}\) for \(\alpha>0\) and \(x_{i}\) of order \(O\left(N^{0}\right)\). One then introduces a density function
\[
\rho(x)=\frac{1}{N} \sum_{i=1}^{N} \delta\left(x-x_{i}\right)
\]
which in the continuum limit should approach an \(L^{1}\) function normalized as
\[
\int d x \rho(x)=1
\]
In terms of this density function, one finds that
\[
F \approx-\frac{9 \pi}{8} N^{2+\alpha} \int d x d y \rho(x) \rho(y)(|x-y|+|x+y|)+\frac{\pi(8-N_{f})}{3} N^{1+3 \alpha} \int d x \rho(x) |x|^{3}
\]
where the large argument expansions have been used, and terms subleading in \(N\) have been dropped. This only has non-trivial saddle points when both terms scale the same with \(N\), which demands that \(\alpha=1 / 2\) and gives the famous result that \(F \propto N^{5 / 2}\). Extremizing the free energy over normalized density functions then gives
\[
F \approx-\frac{9 \sqrt{2} \pi N^{5 / 2}}{5 \sqrt{8-N_{f}}}
\]
This value of the free energy is to be identified with the renormalized on-shell action of the supersymmetric \(\mathrm{AdS}_{6}\) solution. This identification yields the following relation between the six-dimensional Newton's constant \(G_{6}\) and the parameters \(N\) and \(N_{f}\) of the dual SCFT,
\[
G_{6}=\frac{5 \pi \sqrt{8-N_{f}}}{27 \sqrt{2}} N^{-5 / 2}
\]
\end{document}
